\chapter{The Design Space of Memory Consistency Models}
\label{ch:models}

This chapter is the formal core of the paper.
It introduces a single vocabulary
of ordering relations, expresses five consistency models as successively weaker
constraints within it, and establishes the containment hierarchy between them.
The formalism follows the execution-witness and access-graph treatment of
Sarangi \cite[Ch.~9]{sarangi2023nextgen}, with terminology reconciled to that of
Nagarajan et al.\ \cite{nagarajan2020primer} where the two differ.

\section{Executions and Relations}
\label{sec:relations}

\subsection{Memory Events}

An \emph{execution} of a multithreaded program is a set $E$ of memory events.
Each event $e \in E$ is a load $\Ld{a}$ or a store $\St{a}$ to some address $a$,
issued by an identified core, and carrying a value.
Non-memory instructions are
irrelevant to the model and are omitted.

\subsection{The Four Primitive Relations}

\begin{description}
\item[Program order, $\po$.] A per-core total order. $e_1 \po e_2$ if both are
  issued by the same core and $e_1$ precedes $e_2$ in that core's instruction
  stream.
This relation is given by the program and is not subject to negotiation.

\item[Reads-from, $\rf$.] $\St{a} \rf \Ld{a}$ if the load returns the value
  written by that store.
Every load reads from exactly one store (taking the
  initialisation of memory as a store), so $\rf$ is a function on loads.

\item[Coherence order, $\co$.] For each address $a$, a total order on all stores
  to $a$. This is exactly the order guaranteed by the coherence protocol of
  Chapter~3, and its existence is the formal content of the SWMR invariant.

\item[From-read, $\fr$.] A derived relation. $\Ld{a} \fr \St{a}$ if the load reads
  from some store that precedes this store in $\co$; informally, the load happened
  before this store overwrote the value it observed. Formally,
  \[
    \fr \;=\; \rf^{-1} \,;\, \co
  \]
  where $;$ denotes relational composition.
\end{description}

\subsection{The Central Definition}

Every model in this chapter has the same shape.
Each defines a subset of program
order, called \emph{preserved program order} and written $\ppo$, consisting of
those program-order pairs the model promises to maintain globally.
An execution
is permitted by the model exactly when the combined relation is acyclic.

\begin{definition}[Model-permitted execution]
\label{def:permitted}
Let $M$ be a consistency model with preserved program order $\ppo_M$. An
execution $E$ is permitted by $M$ if and only if the relation
\[
  \ppo_M \;\cup\; \rf \;\cup\; \co \;\cup\; \fr
\]
is acyclic over $E$.
\end{definition}

A cycle in this relation is a proof that no total order of events consistent with
the model can produce the execution.
The models therefore differ in exactly one
respect --- how much of $\po$ survives into $\ppo$ --- and the entire design space
can be presented as a single table.
That table is \cref{tab:models}.

\section{Litmus Tests}
\label{sec:litmus-formal}

A \emph{litmus test} is a minimal multithreaded program, together with a
distinguished final state, used to distinguish models.
If a model permits the
state, the hardware may produce it; if not, it may not.

\subsection{SB: Store Buffering}

Listing~\ref{lst:sb} of Chapter~1, in litmus form:
\begin{center}
\begin{tabular}{@{}ll@{}}
\toprule
\textbf{Core 0} & \textbf{Core 1} \\
\midrule
\code{St(x,1)} & \code{St(y,1)} \\
\code{Ld(y)} $\to$ \code{r0} & \code{Ld(x)} $\to$ \code{r1} \\
\bottomrule
\end{tabular}
\quad
\emph{Target:} $\texttt{r0}=0 \wedge \texttt{r1}=0$
\end{center}

\begin{figure}[H]
\centering
\begin{tikzpicture}[
    ev/.style={draw, rounded corners=2pt, minimum width=20mm, minimum height=7mm,
               font=\ttfamily\small, fill=black!4},
    rel/.style={-{Latex[length=2mm]}, thick},
    rl/.style={font=\scriptsize\itshape}]

  \node[ev] (a) at (0,2)   {St(x,1)};
  \node[ev] (b) at (0,0)   {Ld(y)=0};
  \node[ev] (c) at (5,2)   {St(y,1)};
  \node[ev] (d) at (5,0)   {Ld(x)=0};

  \draw[rel] (a) -- node[rl,left]{po} (b);
  \draw[rel] (c) -- node[rl,right]{po} (d);
  \draw[rel] (b) -- node[rl,above,pos=0.5]{fr} (c);
  \draw[rel] (d) -- node[rl,below,pos=0.5]{fr} (a);
\end{tikzpicture}
\caption{Access graph for the SB litmus test with the target outcome. The four
edges form a cycle, so the outcome is forbidden by any model that preserves
store~$\to$~load program order. TSO does not preserve it, and the cycle is
therefore not a cycle in $\ppo_{\mathrm{TSO}} \cup \rf \cup \co \cup \fr$.}
\label{fig:sb-graph}
\end{figure}

Each load reads the initial value, so each load is $\fr$-before the other core's
store.
Together with the two $\po$ edges this closes a cycle.
Under SC, where
$\ppo = \po$, the cycle is present and the outcome is forbidden.
Under TSO, where
store~$\to$~load pairs are dropped from $\ppo$, both $\po$ edges vanish and no
cycle remains: the outcome is permitted.

\subsection{MP: Message Passing}

\begin{center}
\begin{tabular}{@{}ll@{}}
\toprule
\textbf{Core 0} & \textbf{Core 1} \\
\midrule
\code{St(data,1)} & \code{Ld(flag)} $\to$ \code{r0} \\
\code{St(flag,1)} & \code{Ld(data)} $\to$ \code{r1} \\
\bottomrule
\end{tabular}
\quad
\emph{Target:} $\texttt{r0}=1 \wedge \texttt{r1}=0$
\end{center}

This is the idiom underlying every publication protocol: write the payload, then
set a flag; on the reader, test the flag, then read the payload.
The target
outcome --- flag observed set, payload observed stale --- requires either that
Core~0's two stores were reordered (a store~$\to$~store relaxation) or that
Core~1's two loads were reordered (a load~$\to$~load relaxation).
TSO permits
neither, and MP is accordingly forbidden on x86.
Weak models permit both, and MP
is observable on ARM.

\subsection{LB, IRIW, WRC and the Coherence Tests}

Four further tests complete the discriminating set used in this paper.

\begin{description}
\item[LB (Load Buffering).] Each core loads one variable and stores the other;
  the target has each load returning the other core's store.
Requires
  load~$\to$~store reordering.
Forbidden under TSO, permitted under ARM and POWER
  in principle, though rarely observed.
\item[IRIW (Independent Reads of Independent Writes).] Two writer cores and two
  reader cores; the target has the readers disagreeing about the order of the two
  writes.
Distinguishes \emph{multicopy atomic} models, in which a store becomes
  visible to all cores simultaneously, from non-multicopy-atomic ones.
POWER and
  ARMv7 permit it; ARMv8, revised in 2018, does not.
\item[WRC (Write-to-Read Causality).] Three cores in a causal chain; tests
  whether causality is transitive.
\item[CoRR, CoWR, CoRW.] Single-address tests that check coherence itself rather
  than consistency.
Every model forbids their target outcomes; they serve as
  controls, and any hardware or simulator producing them is defective.
\end{description}

\section{The Models}
\label{sec:the-models}

\subsection{Sequential Consistency}

\begin{definition}[Sequential consistency, after Lamport \cite{lamport1979multiprocessor}]
A multiprocessor is sequentially consistent if the result of any execution is the
same as if the operations of all processors were executed in some sequential
order, and the operations of each individual processor appear in that sequence in
the order specified by its program.
\end{definition}

In the framework of \cref{def:permitted}, $\ppo_{\mathrm{SC}} = \po$: all four
program-order pairs are preserved. SC is the model that validates the informal
reasoning of \S\ref{sec:motivating}, and it is the only model in this chapter
that does.

\subsection{Total Store Order}

TSO relaxes exactly one pair: a store may be reordered with a \emph{later} load to
a different address.
This corresponds precisely to a FIFO store buffer with
store-to-load forwarding, and nothing else.
It is the model of SPARC TSO and,
with minor differences of specification, of x86.

Two further properties characterise TSO. The store buffer is FIFO, so
store~$\to$~store order is preserved. And TSO is \emph{multicopy atomic}: a store
that leaves a core's buffer becomes visible to all other cores at once, which is
why IRIW is forbidden.

\subsection{Partial Store Order}

PSO additionally relaxes store~$\to$~store, modelling a non-FIFO store buffer that
may drain entries out of order or coalesce them.
It was defined by SPARC and is
of limited practical importance today, but it is a necessary intermediate point in
the hierarchy and it isolates the store~$\to$~store relaxation for study.

\subsection{Weak Ordering and RMO}

Weak ordering, introduced by Dubois, Scheurich and Briggs
\cite{dubois1986memory}, abandons the attempt to preserve any program order
between ordinary accesses.
It instead divides accesses into \emph{data} and
\emph{synchronisation} operations and guarantees ordering only at synchronisation
points.
SPARC RMO, Alpha, POWER, ARMv7 and RISC-V's RVWMO are all of this family,
differing in detail --- notably in whether they are multicopy atomic and in which
dependencies are respected.

Three dependencies survive in essentially every weak model, because violating them
would break single-threaded semantics:

\begin{description}
\item[Address dependency.] A load whose result computes the address of a later
  access.
\item[Data dependency.] A store whose value is computed from an earlier load.
\item[Control dependency.] A branch conditional on a load. Notably, control
  dependencies order later \emph{stores} but not later \emph{loads}, since the
  latter may be executed speculatively past the branch --- a subtlety that has
  caused documented bugs in production code.
\end{description}

\subsection{Release Consistency}

Release consistency, due to Gharachorloo et al.\ \cite{gharachorloo1990memory},
refines weak ordering by distinguishing two kinds of synchronisation:

\begin{itemize}
\item An \emph{acquire} orders itself before all subsequent accesses.
\item A \emph{release} orders all preceding accesses before itself.
\end{itemize}

Neither constrains operations in the other direction, so the fences are strictly
cheaper than a full barrier.
This is the model exposed by C++11's
\code{memory\_order\_acquire} and \code{memory\_order\_release} and implemented
directly in ARMv8's \code{LDAR} and \code{STLR} instructions.

\subsection{Summary Table}

\begin{table}[H]
\centering
\small
\caption{Program-order pairs preserved by each model. A tick indicates the pair is
in $\ppo$; a cross indicates it may be reordered.}
\label{tab:models}
\begin{tabular}{@{}lcccccl@{}}
\toprule
\textbf{Model} & \textbf{S$\to$L} & \textbf{S$\to$S} & \textbf{L$\to$L} &
\textbf{L$\to$S} & \textbf{MCA} & \textbf{Deployed in} \\
\midrule
SC        & \checkmark & \checkmark & \checkmark & \checkmark & \checkmark & --- \\
TSO       & $\times$   & \checkmark & \checkmark & \checkmark & \checkmark & x86, SPARC \\
PSO       & $\times$   & $\times$   & \checkmark & \checkmark & \checkmark & SPARC \\
Weak/RMO  & $\times$   & $\times$   & $\times$   & $\times$   & varies     & ARM, POWER, RISC-V \\
RC        & $\times$   & $\times$   & $\times$   & $\times$   & varies     & C++11, ARMv8 \\
\bottomrule
\end{tabular}
\end{table}

\noindent MCA denotes multicopy atomicity. ARMv8 became multicopy atomic in its
2018 revision; POWER and ARMv7 are not.

\section{The Containment Hierarchy}
\label{sec:containment}

Let $\mathcal{X}(M)$ denote the set of executions permitted by model $M$.

\begin{theorem}[Containment]
\label{thm:containment}
$\mathcal{X}(\mathrm{SC}) \subsetneq \mathcal{X}(\mathrm{TSO}) \subsetneq
 \mathcal{X}(\mathrm{PSO}) \subsetneq \mathcal{X}(\mathrm{Weak})$.
\end{theorem}

\begin{proof}[Proof sketch]
\emph{Containment.} By construction $\ppo_{\mathrm{Weak}} \subseteq
\ppo_{\mathrm{PSO}} \subseteq \ppo_{\mathrm{TSO}} \subseteq \ppo_{\mathrm{SC}}$,
since each model is obtained from its predecessor by deleting pairs.
If a relation
is acyclic then any subrelation of it is acyclic; hence a superset of $\ppo$
generates a superset of the cycles, and so a subset of the permitted executions.

\emph{Strictness.} Each inclusion is witnessed by a litmus test permitted by the
weaker model and forbidden by the stronger: SB separates SC from TSO; MP with the
stores reordered separates TSO from PSO; MP with the loads reordered separates PSO
from Weak.
\Cref{fig:sb-graph} gives the argument in the first case; the others
are structurally identical and are given in full in Appendix~B.
\end{proof}

The hierarchy has a practical reading.
A program correct under a weak model is
correct under every stronger one, so code written for ARM runs correctly on x86.
The converse fails, and this asymmetry is the source of a well-known class of
porting defect.

\section{Fences}
\label{sec:fences}

A fence restores ordering that the model does not provide by default.
Its
semantics are stated as an addition to $\ppo$: a full fence $F$ contributes
$\{(e_1,e_2) : e_1 \po F \po e_2\}$.

\begin{table}[H]
\centering
\small
\caption{Ordering instructions in deployed instruction sets.}
\label{tab:fences}
\begin{tabular}{@{}llp{6.3cm}@{}}
\toprule
\textbf{ISA} & \textbf{Instruction} & \textbf{Effect} \\
\midrule
x86    & \code{MFENCE}         & Full barrier; drains the store buffer. \\
x86    & \code{LOCK} prefix    & Atomic RMW; carries full-barrier semantics. \\
ARMv8  & \code{DMB ISH}        & Full data memory barrier, inner shareable. \\
ARMv8  & \code{LDAR} / \code{STLR} & Load-acquire, store-release. \\
POWER  & \code{sync} / \code{lwsync} & Heavyweight and lightweight sync. \\
RISC-V & \code{FENCE} \emph{pred,succ} & Bit-encoded predecessor/successor sets. \\
\bottomrule
\end{tabular}
\end{table}

The RISC-V encoding is worth noting: rather than offering a fixed menu, it exposes
four predecessor and four successor bits, so that the programmer requests exactly
the orderings needed. This is the most explicit acknowledgement in any deployed
ISA that fence strength is a continuum rather than a binary.

Fences are not free.
Restoring an ordering means reintroducing the stall the
relaxation was introduced to avoid; a full barrier on x86 costs on the order of
$50$--$100$ cycles.
Quantifying this cost on real hardware is one of the
measurements specified in Chapter~6.

\section{Data-Race-Free Programming}
\label{sec:drf}

The models above are hardware contracts.
Programmers work in languages, and the
bridge between the two is the data-race-freedom theorem of Adve and Hill
\cite{adve1990weak}.

\begin{definition}[Data race]
Two accesses conflict if they target the same address and at least one is a store.
An execution has a data race if two conflicting accesses from different cores are
unordered by synchronisation.
A program is data-race-free (DRF) if no sequentially
consistent execution of it contains a race.
\end{definition}

\begin{theorem}[DRF, informally]
A DRF program executing on a weakly ordered machine that respects
synchronisation operations produces only sequentially consistent results.
\end{theorem}

This is the settlement on which modern practice rests, and it is the intellectual
foundation of both the Java Memory Model \cite{manson2005java} and the C++11
model \cite{boehm2008foundations}.
Its bargain is explicit: \emph{annotate your
synchronisation correctly and you may reason sequentially}.
Its weakness is
equally explicit --- the guarantee is void for racy programs, and racy programs
are common.
Chapter~11 returns to whether this bargain was a good one.

\section{Summary}

Five models, one formalism, one table.
The models differ only in how much of
program order they preserve, and they form a strict hierarchy.
What the table does
not show is cost: it says what each model forbids, not what forbidding it is worth.
Chapter~5 recounts how the field came to settle where it did, and Chapter~6 sets
out how this paper proposes to measure whether that settlement was correct.
